% Anexos (cargado por main.tex tras \appendix)

% Glosario movido a frontmatter/8_glossary.tex (2026-07-10, punto 1 de la tutora).

\chapter{Conjunto de características}
\label{anexo:features}

El especialista final emplea 36 características causales sin variables de reloj
(\emph{no\_clock}), agrupadas por bloques:

\begin{itemize}[noitemsep]
\item \textbf{Estado local del libro de Polymarket:} precio medio, \emph{spread}, microprice,
desbalance de órdenes, coste visible de entrada y deltas recientes del precio.
\item \textbf{Referencias externas:} retornos orientados del mercado al contado y de los
futuros perpetuos a 2 y 8 segundos; medidas de consenso y de actualización.
\item \textbf{Volatilidad:} volatilidad realizada del perpetuo a 5 segundos, usada en el
diagnóstico posterior de cambio de régimen.
\item \textbf{Contexto de ventana:} fase de la ventana \emph{prestart} y proxies de calidad del
libro (cobertura, indicadores de degradación).
\end{itemize}

Se excluyen deliberadamente las variables de reloj (hora, minuto, segundo) para evitar que el
modelo aprenda atajos temporales triviales.

\chapter{Hiperparámetros de los modelos}
\label{anexo:hiperparametros}

\begin{table}[h]
\centering
\caption{Hiperparámetros principales del especialista (HistGradientBoosting, dos cabezas).}
\begin{tabular}{ll}
\toprule
\textbf{Hiperparámetro} & \textbf{Valor} \\
\midrule
Tasa de aprendizaje & 0,045 \\
Iteraciones máximas & 220 (con parada anticipada) \\
Fracción de validación interna & 0,15 \\
Máximo de hojas por árbol & 15 \\
Mínimo de muestras por hoja & 120 \\
Regularización L2 & 0,10 \\
\bottomrule
\end{tabular}
\end{table}

\begin{table}[h]
\centering
\caption{Hiperparámetros principales de los modelos profundos (línea de exploración).}
\begin{tabular}{ll}
\toprule
\textbf{Hiperparámetro} & \textbf{Valor} \\
\midrule
Optimizador & AdamW \\
Tasa de aprendizaje & 0,0015 \\
Decaimiento de pesos & $10^{-4}$ \\
Tamaño de lote & 192 \\
Paciencia (parada anticipada) & 18 épocas \\
Dimensión oculta de la GRU & 56 \\
\emph{Dropout} (\emph{BookEncoder}) & 0,08 \\
\bottomrule
\end{tabular}
\end{table}

El especialista base congelado opera con dos umbrales: valor esperado predicho superior a
0,75 y probabilidad de salud no inferior a 0,50. Tras el diagnóstico del bloque fuera de
muestra se propone una política futura que exige, además, volatilidad realizada a 5 segundos
no superior a 0,6657 (percentil 80 del entrenamiento). La incertidumbre exploratoria se
estima con 5\,000 remuestreos y semilla fija, tanto por filas como agrupando por mercado.

\chapter{Pre-registro v2 y verificación del sello}
\label{anexo:preregistro}

Este anexo reproduce el compromiso de pre-registro que sostiene la validación
prospectiva del Capítulo~\ref{ch::resultados} y detalla cómo verificarlo de forma
independiente. El documento íntegro vive en el repositorio
(\texttt{docs/PREREGISTRO\_V2.md}); aquí se resume su contenido comprobable.

\section*{Política congelada}

La señal y la política se fijaron en el entrenamiento de mayo y no se han modificado
desde entonces: universo \emph{prestart} estricto ($-60$ a $-45$\,s respecto a la
apertura, \emph{spread} $\le 2$ ticks, libro saludable, coste visible acotado), dos
cabezas de \emph{gradient boosting} sobre el conjunto de características \emph{no\_clock},
y regla de operación «valor esperado $>0{,}75$ y salud $\ge 0{,}5$». La integridad de
los artefactos se ancla con sus resúmenes SHA-256, incluidos literalmente en el
pre-registro (modelos de valor esperado y de salud, manifiesto de características y
\emph{gate} GARCH congelado).

\section*{Brazos de \emph{gate} y capas de reporte}

Se comprometieron por adelantado tres brazos de \emph{gate} —estático heredado,
GARCH(1,1) ajustado solo con datos anteriores a junio, y clustering de regímenes
(adenda A)— y tres capas de reporte que se publican todas: net@0,5 histórica, taker con
comisión oficial (con expectativa \emph{declarada como negativa} antes de mirar los
datos) y maker simulado. Declarar la expectativa negativa por adelantado es lo que
convierte el resultado taker en una confirmación y no en una racionalización
\emph{a posteriori}.

\section*{Cortes prospectivos}

Los datos elegibles son la captura de la Pi posterior a la reanudación del recolector,
procesados con el mismo \emph{pipeline} congelado. Se fijaron dos cortes —el primero a
finales de agosto, el segundo justo antes de la entrega— sin ningún vistazo intermedio a
las métricas de resultado.

\section*{Verificación independiente del sello temporal}

El compromiso se ancla en la cadena de Bitcoin mediante
OpenTimestamps~\cite{opentimestamps}. Para verificarlo sin confiar en el autor:

\begin{enumerate}[noitemsep]
\item Obtener el \emph{commit} del repositorio que introduce el pre-registro y su prueba
\texttt{.ots} asociada (\texttt{docs/prereg\_v2\_commit.ots}).
\item Ejecutar \texttt{ots verify docs/prereg\_v2\_commit.ots}: la herramienta consulta
un nodo de Bitcoin y confirma el bloque —y por tanto la fecha— en que quedó atestado el
resumen del documento.
\item Comprobar que esa fecha es anterior al primer dato de la ventana de evaluación
prospectiva. Como la atestación de Bitcoin no puede retro-fecharse, la anterioridad del
compromiso queda demostrada sin necesidad de confiar en la palabra del autor ni en un
sello privado.
\end{enumerate}

Una fe de erratas posterior (la corrección del modelo de comisiones,
Sección~\ref{sec::fe_erratas_fees}) se selló con el mismo procedimiento, de modo que la
cadena de compromisos es completa y auditable.

\chapter{Auditoría del ledger público}
\label{anexo:ledger}

Las cifras económicas centrales del trabajo pueden recalcularse de forma independiente sin
acceso al corpus privado ni a los modelos entrenados, a partir de un único fichero
anonimizado que se publica con el código: el ledger de acciones candidatas
(\texttt{results/\allowbreak final\_candidate\_\allowbreak actions\_\allowbreak anonymized.csv}). Este anexo describe su
contenido y el procedimiento de auditoría.

\section*{Contenido del ledger}

El ledger contiene una fila por unidad de acción del bloque fuera de muestra (754 filas).
Cada fila registra un identificador de acción, el día de sesión, los identificadores de
mercado y de contrato \emph{anonimizados} mediante \emph{hash} (de modo que no se pueda
reconstruir la identidad del mercado), el instante relativo a la apertura, el régimen de
volatilidad (\texttt{high}/\texttt{low}) y el resultado neto en ticks a tres niveles de
coste de ejecución (medio tick de referencia, y las variantes de un cuarto y de un tick
completo). No incluye características del modelo ni datos crudos: es exactamente la
información necesaria para auditar el resultado, y nada más.

\section*{Procedimiento de reproducción}

Con solo ese fichero y una utilidad de análisis de datos, la auditoría reproduce las tres
cifras que sostienen las conclusiones económicas:

\begin{enumerate}[noitemsep]
\item \textbf{Resultado confirmatorio.} La media del resultado neto al coste de referencia
sobre las 754 acciones (241 mercados, cinco días) es de $+0{,}349$ ticks por acción,
idéntica a la cifra publicada. Bajo el coste conservador de un tick completo el promedio
cae a $-0{,}026$: la señal no sobrevive al coste de estrés.
\item \textbf{Diagnóstico \emph{post-hoc} de baja volatilidad.} Restringiendo a las 318
acciones de régimen de baja volatilidad (156 mercados), la media asciende a $+1{,}069$
ticks. Se recuerda que esta separación se decidió tras inspeccionar el bloque, por lo que
es una hipótesis, no una segunda validación.
\item \textbf{Incertidumbre agrupada por mercado.} Un \emph{bootstrap} que remuestrea
mercados enteros —para no tratar como independientes las hasta ocho señales solapadas de un
mismo mercado— sitúa el intervalo de confianza del 90\,\% del corte de baja volatilidad en
$[+0{,}193,\,+2{,}004]$ ticks: positivo pero ancho, coherente con el escaso soporte
temporal.
\end{enumerate}

El cuaderno \texttt{08\_auditoria\_final\_ledger} ejecuta estos tres cálculos paso a paso y
constituye la vía pública más directa para verificar el resultado final del trabajo con
independencia del autor.
